\documentclass[11pt]{article}

\usepackage[utf8]{inputenc}
\usepackage[T1]{fontenc}
\usepackage[spanish,es-noquoting]{babel}
\usepackage{lmodern}
\usepackage{amsmath,amssymb,amsfonts}
\usepackage[a4paper,margin=2.3cm]{geometry}
\usepackage{booktabs}
\usepackage{array}
\usepackage{enumitem}
\usepackage{xcolor}
\usepackage{tcolorbox}
\usepackage{titlesec}
\usepackage{parskip}
\usepackage{hyperref}

\hypersetup{colorlinks=true,linkcolor=blue!55!black,urlcolor=blue!55!black}
\definecolor{boxblue}{RGB}{225,235,248}
\definecolor{rulegray}{RGB}{120,120,120}

\titleformat{\section}{\Large\bfseries\color{blue!40!black}}{\thesection}{0.6em}{}
\titleformat{\subsection}{\large\bfseries}{\thesubsection}{0.6em}{}

\newtcolorbox{keybox}{colback=boxblue,colframe=blue!40!black,boxrule=0.4pt,
  arc=2pt,left=6pt,right=6pt,top=4pt,bottom=4pt}

\newcommand{\E}{\mathbb{E}}
\newcommand{\R}{\mathbb{R}}
\newcommand{\Var}{\operatorname{Var}}
\newcommand{\Cov}{\operatorname{Cov}}
\newcommand{\Om}{\Omega_n}
\newcommand{\sg}{\sigma}

\title{\textbf{Resumen para el examen}\\[2pt]
\large Estadística geométrica: $L^2$, proyección y esperanza condicional}
\author{}
\date{}

\begin{document}
\maketitle
\vspace{-2.5em}


% ============================================================
\section*{0. Marco teórico común (la llave de todo)}
\addcontentsline{toc}{section}{0. Marco teórico común}

\subsection*{0.1 Espacio empírico de probabilidad}
Una base de datos es la terna $(\Om,\mathcal P(\Om),P_n)$ donde:
\begin{itemize}[leftmargin=1.4em,itemsep=1pt]
\item $\Om=\{\omega_1,\dots,\omega_n\}$ es el \textbf{conjunto de etiquetas (labels)} = los
\textbf{individuos estadísticos} (identifican cada fila de forma única). No confundir la etiqueta
con el valor de una variable.
\item Probabilidad \textbf{uniforme} (Laplace): $p_\omega=1/n$, y $P_n(A)=|A|/n$.
\item Una \textbf{variable} es una función $X:\Om\to\mathrm{Im}(X)$. Cada columna es una variable.
\end{itemize}

\subsection*{0.2 Partición inducida por una variable}
Toda $X$ con $\mathrm{Im}(X)=\{x_1,\dots,x_r\}$ induce una \textbf{partición} de $\Om$ en sus
\textbf{fibras / bloques}: $A_i:=X^{-1}(\{x_i\})=\{\omega:X(\omega)=x_i\}$. No vacíos, disjuntos,
cubren $\Om$. $\sg(X)$ designa el \textbf{subespacio de funciones constantes sobre cada bloque}.

\subsection*{0.3 $L^2$ como espacio de Hilbert}
$L^2(\Om,\mathcal P(\Om),P_n)$ con
\[
\langle U,V\rangle=\E[UV]=\sum_\omega U(\omega)V(\omega)P_n(\{\omega\})=\tfrac1n\sum_{i=1}^n U(\omega_i)V(\omega_i),
\qquad \|U\|^2=\E[U^2].
\]
Como $\Om$ es finito, toda función está en $L^2$ y el espacio es completo (Hilbert).

\subsection*{0.4 Esperanza condicional = proyección ortogonal}
\begin{keybox}
\[ E(Y\mid\sg(X))=P_{\sg(X)}Y=\text{proyección ortogonal de }Y\text{ sobre }\sg(X). \]
\end{keybox}
Es el único $\hat Y\in\sg(X)$ que minimiza $\|Y-\hat Y\|^2$ (mejor predicción en media cuadrática).
En cada individuo vale \textbf{el promedio de $Y$ sobre su bloque}:
\[ E(Y\mid\sg(X))(\omega)=\frac1{|B|}\sum_{\omega'\in B}Y(\omega'),\quad B=\text{bloque de }\omega. \]
\textbf{Teorema de la Proyección (hecho central):} el residuo $Z=Y-\hat Y$ es ortogonal a todo el
subespacio: $\langle Z,W\rangle=0$ para todo $W\in\sg(X)$.\\
\textbf{Dato clave:} la función constante $\mathbf 1$ siempre está en $\sg(X)$.

\subsection*{0.5 Tres consecuencias inmediatas}
\begin{enumerate}[leftmargin=1.6em,itemsep=2pt]
\item \textbf{Insesgamiento / esperanza total:} como $\mathbf 1\in\sg(X)$ y $Z\perp\sg(X)$,
$\E[Z]=\langle Z,\mathbf 1\rangle=0\Rightarrow \E[E(Y\mid\sg(X))]=\E[Y]$.
\item \textbf{Pitágoras (norma):} $\|Y\|^2=\|\hat Y\|^2+\|Z\|^2$.
\item \textbf{Descomposición de varianza (ANOVA):}
\[ \Var(Y)=\underbrace{\Var(E(Y\mid\sg(X)))}_{\text{entre bloques (SC}_\text{inter})}
+\underbrace{\E(\Var(Y\mid\sg(X)))}_{\text{dentro (SC}_\text{intra})}. \]
\end{enumerate}

\subsection*{0.6 Fórmulas de proyección útiles}
\begin{itemize}[leftmargin=1.4em,itemsep=1pt]
\item Sobre una recta generada por $v\neq0$: $P_{\sg(v)}Y=\dfrac{\langle Y,v\rangle}{\langle v,v\rangle}v$.
\item \textbf{Suma de proyecciones SOLO si hay ortogonalidad:} si $\mathcal M\perp\mathcal N$ entonces
$P_{\mathcal M\oplus\mathcal N}=P_{\mathcal M}+P_{\mathcal N}$. Si no, \textbf{falla} (ecuaciones normales).
\end{itemize}

% ============================================================
\section*{1. Base de datos (reformas fiscales)}
\textbf{Planteo:} tabla País$\times$Año (20 individuos, 9 variables). Espacio $\R^{20}$, pesos $1/20$.

\begin{enumerate}[leftmargin=1.6em,itemsep=2pt]
\item \textbf{Etiquetas:} los pares $(\text{País},\text{Año})$: $I=\{(\mathrm{FRA},2015),\dots\}$,
$|I|=4\times5=20$. (El enunciado nombra 5 países pero la tabla trae 4.)
\item \textbf{Partición por FR:} binaria $\Rightarrow$ 2 clases: $A_0=\{FR=0\}$ (2015--16, $|A_0|=8$)
y $A_1=\{FR=1\}$ (2017--19, $|A_1|=12$).
\item \textbf{$E(\text{Inv}\mid\text{Año})$:} 5 clases (4 individuos c/u), valor = promedio del año:
2015$\to$18.175, 2016$\to$17.300, 2017$\to$18.825, 2018$\to$21.000, 2019$\to$18.325.
Media global $\bar x=374.5/20=18.725$.
\item \textbf{$r^2$ e interpretación:} razón de correlación $\eta^2$.
\end{enumerate}

\textbf{Cálculo del $r^2$:}
\[ SC_\text{inter}=\sum_a n_a(\bar x_a-\bar x)^2=30.715,\quad
SC_\text{total}=\sum_i x_i^2-20\bar x^2=53.8975,\quad SC_\text{intra}=23.1825. \]
\[ r^2=\frac{SC_\text{inter}}{SC_\text{total}}=\frac{30.715}{53.8975}\approx0.570\ (\approx57\%). \]
\textbf{Interpretación geométrica:} $E(\text{Inv}\mid\text{Año})$ es la proyección ortogonal de Inv
sobre el subespacio $F$ de funciones constantes por año; el residuo es ortogonal a $F$. Con vectores
centrados, $\|\text{Inv}-\bar x\|^2=\|\hat Y-\bar x\|^2+\|\text{Inv}-\hat Y\|^2$ y $r^2=\cos^2\theta$,
con $\cos\theta=\sqrt{0.57}\approx0.755\Rightarrow\theta\approx41^\circ$.

% ============================================================
\section*{2. Coherencia en una base de datos (falacia ecológica)}
\textbf{Planteo:} votos identificados por folio único; cada folio tiene opción de voto y pertenece a
mesa $\to$ local $\to$ distrito.
\begin{enumerate}[leftmargin=1.6em,itemsep=2pt]
\item \textbf{Etiquetas = los folios.} La opción de voto NO es etiqueta, es una variable
$V:\Omega_N\to\{\text{cand.1, cand.2, nulo, blanco}\}$. Los folios deben ser únicos, exhaustivos y
mutuamente excluyentes.
\item \textbf{Mesa, Local, Distrito SÍ son variables}, con particiones \textbf{jerárquicas/encajonadas}:
$\sg(D)\subseteq\sg(L)\subseteq\sg(M)\subseteq L^2(\Omega_N)$ (Mesa más fina, Distrito más gruesa).
\item \textbf{Tasa de cesantía comunal $C$:} (i) \emph{coherencia jerárquica:} asignación unívoca
folio$\to$comuna, $K:\Omega_N\to\{\text{comunas}\}$; (ii) \emph{medibilidad:} $C$ constante por comuna,
i.e. $C\in\sg(K)$.
\end{enumerate}
\textbf{El punto de fondo (falacia ecológica):} como $C\in\sg(K)$ es constante dentro de cada comuna,
solo aporta variabilidad \emph{entre} comunas. Toda regresión factoriza por $E(V\mid C)=E(V\mid\sg(K))$:
\[ \Var(V)=\underbrace{\Var(E(V\mid\sg(K)))}_{\text{entre comunas (ve }C)}
+\underbrace{\E(\Var(V\mid\sg(K)))}_{\text{dentro: invisible}}. \]
Atribuir el efecto agregado al individuo es la \textbf{falacia ecológica} (Robinson, 1950); puede
invertirse (\textbf{paradoja de Simpson}). Para el nivel individual hace falta una variable definida
sobre el folio con partición más fina que la comunal.

% ============================================================
\section*{3. Modelos lineales encajonados (nested)}
\textbf{Modelos:} M1: $E(Y\mid X_1)=\alpha_0+\alpha_1X_1$; M2: $E(Y\mid X_1,X_2)=\alpha_0+\alpha_1X_1+\alpha_2X_2$.\\
\textbf{Marco:} $\Omega_3$, uniforme $1/3$, $\langle U,V\rangle=\tfrac13\sum U_iV_i$.

\textbf{Ejemplo 1 — regresores ORTOGONALES (se cumple la aditividad).}
$\mathbf 1=(1,1,1)$, $X=(-1,0,1)$. Como $\langle\mathbf 1,X\rangle=0$:
\[ E(Y\mid\sg(\mathbf 1,X))=\bar y\,\mathbf 1+\tfrac{y_3-y_1}{2}X=E(Y\mid\sg(\mathbf 1))+E(Y\mid\sg(X)).\ \checkmark \]
La igualdad se debe a la \textbf{ortogonalidad}: proyectar sobre el plano = sumar proyecciones sobre cada eje (Pitágoras).

\textbf{Ejemplo 2 — regresores NO ortogonales (falla).}
$\mathbf 1=(1,1,1)$, $Z=(1,0,1)$, $\langle\mathbf 1,Z\rangle=\tfrac23\neq0$. Por \textbf{ecuaciones normales}:
\[ a=y_2,\quad b=\tfrac{y_1+y_3}{2}-y_2,\quad E(Y\mid\sg(\mathbf 1,Z))=\Big(\tfrac{y_1+y_3}{2},y_2,\tfrac{y_1+y_3}{2}\Big). \]
En $\omega_2$: $y_2$ vs $\tfrac{y_1+y_2+y_3}{3}$; coinciden solo si $2y_2=y_1+y_3$. En general la suma de
proyecciones cuenta dos veces la parte común.

\textbf{Respuesta (¿cuándo son ``encajonados''?):}
\begin{enumerate}[leftmargin=1.6em,itemsep=2pt]
\item \textbf{Inclusión de subespacios:} $\sg(\mathbf 1)\subseteq\sg(\mathbf 1,X_1)\subseteq\sg(\mathbf 1,X_1,X_2)$
\emph{siempre}. Solo valida comparar ajustes (M2 nunca ajusta peor).
\item \textbf{Descomposición aditiva:} se cumple \emph{si y solo si} $X_2\perp\sg(\mathbf 1,X_1)$.
\end{enumerate}
\begin{keybox}
\textbf{Criterio:} M1 y M2 son ``encajonados'' en sentido fuerte solo cuando $X_2\perp\sg(\mathbf 1,X_1)$.
Si los predictores están correlacionados, agregar $X_2$ altera el coeficiente de $X_1$. La forma correcta es
\textbf{ortogonalizar primero} (\textbf{teorema de Frisch--Waugh--Lovell}). La receta ingenua ``ajusto,
agrego, comparo'' solo es interpretable bajo ortogonalidad.
\end{keybox}

% ============================================================
\section*{4. Teorema de Pitágoras ($R^2$ vs $r_n^2$)}
\textbf{Setup:} $Y=\hat Y+Z$, $\hat Y=E(Y\mid\sg(X))$, $Z=Y-\hat Y$, $\|Y\|^2=\|\hat Y\|^2+\|Z\|^2$.

\textbf{Parte 1 — $\E(Y-E(Y\mid\sg(X)))=0$:} como $\mathbf 1\in\sg(X)$ y $Z\perp\sg(X)$,
$\langle Z,\mathbf 1\rangle=\E[Z]=0$. Consecuencia: $\E[\hat Y]=\E[Y]=:\mu$.

\textbf{Parte 2 — equivalencia (1)$\Leftrightarrow$(2):} con $\Var(U)=\E[U^2]-[\E(U)]^2$ y
$\E[\hat Y]=\E[Y]=\mu$, $\E[Z]=0$. Restando $\mu^2$ a $\E[Y^2]=\E[\hat Y^2]+\E[Z^2]$:
\[ \Var(Y)=\Var(E(Y\mid\sg(X)))+\Var(Y-E(Y\mid\sg(X))). \]
``Restar $\mu^2$'' es reversible $\Rightarrow$ (1) y (2) equivalentes.

\textbf{Parte 3 — comparar.} Con $a=\|\hat Y\|^2=\E[\hat Y^2]$, $b=\|Y\|^2=\E[Y^2]$:
\[ r_n^2=\frac{\|E(Y\mid\sg(X))\|^2}{\|Y\|^2}=\frac{a}{b}\ \text{(crudos)},\qquad
R^2=\frac{\Var(\hat Y)}{\Var(Y)}=\frac{a-\mu^2}{b-\mu^2}\ \text{(varianzas)}. \]
Como $b=a+\E[Z^2]\Rightarrow a\le b$, y por Jensen $a\ge\mu^2$: $\mu^2\le a\le b$. Producto cruzado:
\begin{keybox}\[ R^2\le r_n^2. \]\end{keybox}
\textbf{Igualdad} sii $\mu=\E[Y]=0$ (centrada) o $a=b$ ($Y=\hat Y$ c.s.). $R^2$ es el genuino coeficiente
de determinación (invariante a traslaciones); $r_n^2$ usa momentos crudos y \textbf{sobrestima}.

\textbf{Ejemplo:} $\Omega=\{1,2,3,4\}$ uniforme, bloques $\{1,2\},\{3,4\}$, $Y=(0,2,2,4)$, $\hat Y=(1,1,3,3)$,
$\mu=2$, $b=6$, $a=5$, $\Var(Y)=2$, $\Var(\hat Y)=1$. Entonces
$r_n^2=\tfrac56\approx0.833$, $R^2=\tfrac12=0.5$ ($R^2<r_n^2$).

% ============================================================
\section*{5. Esperanza condicional (construcción y ley total)}
\textbf{Parte 1 — $P:H\to M$ es proyección ortogonal} ($M$ cerrado):
\begin{itemize}[leftmargin=1.4em,itemsep=1pt]
\item \textbf{Definición:} $Px:=y_0$, el único minimizador de $\|x-y\|$ sobre $M$.
\item \textbf{Lema clave:} $y_0$ minimiza $\iff(x-y_0)\perp M$. ($\Leftarrow$) Pitágoras; ($\Rightarrow$)
$\varphi(t)=\|x-(y_0+tm)\|^2$ tiene mínimo en $t=0\Rightarrow\langle x-y_0,m\rangle=0$.
\item \textbf{Lineal} (por unicidad), \textbf{idempotente} ($P^2=P$), \textbf{autoadjunto}
($\langle Px,z\rangle=\langle x,Pz\rangle$), \textbf{contracción} ($\|P\|\le1$), $H=M\oplus M^\perp$.
\end{itemize}
\textbf{Parte 2 — la proyección ES la esperanza condicional:} en $L^2(\Omega,\mathcal F,P)$ con
$\langle X,Y\rangle=\E[XY]$, $M=L^2(\Omega,\mathcal G,P)$ es cerrado. La definición de $E[X\mid\mathcal G]$
($\mathcal G$-medible con $\E[(X-Z)W]=0\ \forall W\in M$) es exactamente la condición de ortogonalidad
$\Rightarrow E[X\mid\mathcal G]=PX$ (mejor predicción en media cuadrática).\\
\textbf{Parte 3 — $\E[E(Y\mid\sg(X))]=\E(Y)$:} como $\mathbf 1\in M$ y $Y-\hat Y\perp M$:
\[ \E[Y]-\E[\hat Y]=\langle Y-\hat Y,\mathbf 1\rangle=0. \]
Verificación: $\hat Y|_{B_j}=\E[Y\mid B_j]$ y $\E[\hat Y]=\sum_j P(B_j)\E[Y\mid B_j]=\E[Y]$ (ley de la esperanza total).

% ============================================================
\section*{6. Igualdad del paralelogramo}
\textbf{Enunciado:} en un espacio con producto interno, $\|x+y\|^2+\|x-y\|^2=2\|x\|^2+2\|y\|^2$.

\textbf{1. Demostración:} expandir con $\|v\|^2=\langle v,v\rangle$, bilinealidad y simetría:
$\|x\pm y\|^2=\|x\|^2\pm2\langle x,y\rangle+\|y\|^2$. Al sumar, los cruzados $\pm2\langle x,y\rangle$
\textbf{se cancelan}.

\textbf{2. Relevancia para la esperanza condicional:} es una \textbf{proyección ortogonal en $L^2$}, cuya
existencia/unicidad (Teorema de la Proyección) usa la igualdad: para una sucesión minimizante $(Y_n)$,
\[ \|Y_m-Y_n\|^2=2\|X-Y_n\|^2+2\|X-Y_m\|^2-4\Big\|X-\tfrac{Y_n+Y_m}{2}\Big\|^2\le 2\|X-Y_n\|^2+2\|X-Y_m\|^2-4d^2\to0, \]
luego $(Y_n)$ es de Cauchy; por completitud converge a $Y^*\in M$. Sin producto interno (norma sin la
igualdad) el argumento se rompe.

\textbf{3. $\|\cdot\|_\infty$ en $\R^n$ ($n\ge2$):} $x=e_1$, $y=e_2$. $\|x\|_\infty=\|y\|_\infty=1$,
$\|x\pm y\|_\infty=1$. LI$=1+1=2$, LD$=2+2=4$. $2\neq4$ $\Rightarrow$ falla.

\textbf{4. $\|\cdot\|_\infty$ en $C[0,1]$:} $F(x)=1$, $G(x)=x$. $\|F\|_\infty=\|G\|_\infty=1$,
$\|F+G\|_\infty=2$, $\|F-G\|_\infty=1$. LI$=4+1=5$, LD$=4$. $5\neq4$ $\Rightarrow$ falla.

% ============================================================
\section*{7. El término de error — dos miradas}
\textbf{Modelo:} $y=X\beta+\varepsilon$, $\mathrm{rango}(X)=p$. Ambas llegan al mismo
$\hat\beta=(X^\top X)^{-1}X^\top y$ por caminos opuestos.

\textbf{Mirada de Gauss (probabilística).} $\varepsilon$ es objeto \textbf{primitivo} con ley supuesta:
$\E[\varepsilon]=0$, $\Var(\varepsilon)=\sigma^2 I_n$ (homocedasticidad + no correlación); para inferir
$\varepsilon\sim\mathcal N(0,\sigma^2I_n)$. La normal surge de exigir que la media sea máximo-verosímil
$\Rightarrow$ MV $=$ minimizar $\sum\varepsilon_i^2$. Sin normalidad, Gauss--Markov: $\hat\beta$ es
\textbf{BLUE}. Estilo: paramétrico, inferencial ($\Var(\hat\beta)=\sigma^2(X^\top X)^{-1}$, tests $t$,$F$, IC).

\textbf{Mirada de la Proyección (geométrica).} $y=\hat y+e$, $\hat y\in\mathcal C(X)$, $e\perp\mathcal C(X)$.
El error \textbf{se define} como complemento ortogonal. Ecuaciones normales $=$ ortogonalidad:
$X^\top e=0$. $P=X(X^\top X)^{-1}X^\top$, $\hat y=Py$, $e=(I-P)y$. Pitágoras $=$ ANOVA sin distribución:
$\|y\|^2=\|\hat y\|^2+\|e\|^2$. En $L^2$, la no correlación del error con los regresores es \textbf{teorema},
no supuesto. Estilo: geométrico, libre de distribución.

\begin{center}
\begin{tabular}{@{}p{4.0cm}p{4.4cm}p{4.4cm}@{}}
\toprule
\textbf{Criterio} & \textbf{Gauss} & \textbf{Proyección}\\
\midrule
Naturaleza del error & Primitivo (v.a. con ley) & Derivado (residuo ortogonal)\\
No correlación & Supuesto & Consecuencia automática\\
Ecuaciones normales & Anular $\partial S/\partial\beta$ & Traducir $X^\top e=0$\\
Por qué MCO & Verosimilitud normal / BLUE & Proyección $=$ punto más cercano\\
Distribución & Central desde el inicio & Postergada (solo para inferir)\\
Estilo & Inferencial/paramétrico & Geométrico/estructural\\
\bottomrule
\end{tabular}
\end{center}
\textbf{Síntesis:} Gauss responde \emph{por qué es óptimo} (inferencial); la proyección, \emph{por qué tiene esa
forma} (estructural). Distinción fina: errores verdaderos $\varepsilon$ $\neq$ residuos $e=(I-P)\varepsilon$.

% ============================================================
\section*{8. Independencia de funciones ($n=13$ primo)}
\textbf{Setup:} $\Om$ empírico uniforme, $X:\Omega_{13}\to\{x_1,\dots,x_r\}$, $Y:\Omega_{13}\to\{y_1,\dots,y_q\}$.

\textbf{1. ¿$r$ arbitrario?} No: $1\le r\le n$, pues $n=\sum_{i=1}^r|A_i|\ge r$ (bloques no vacíos).
$r=1\iff X$ constante; $r=n\iff X$ inyectiva. Con $n=13$: $r\in\{1,\dots,13\}$.

\textbf{2. Partición de $(X,Y)$:} la fibra es $(X,Y)^{-1}(\{(x_i,y_j)\})=A_i\cap B_j$, luego la partición
es $\{A_i\cap B_j:A_i\cap B_j\neq\varnothing\}$: el \textbf{refinamiento común} de $\{A_i\}$ y $\{B_j\}$
(más bloques $\le r\cdot q$ y $\le n$).

\textbf{3. ¿Qué $Y$ son independientes de $X$ fija?} Independencia $\iff
P(X=x_i,Y=y_j)=P(X=x_i)P(Y=y_j)$. En el espacio empírico:
\[ \frac{|A_i\cap B_j|}{13}=\frac{|A_i|}{13}\cdot\frac{|B_j|}{13}\iff |A_i\cap B_j|=\frac{|A_i||B_j|}{13}. \]
El lado izquierdo es entero $\Rightarrow 13\mid|A_i||B_j|$. Como \textbf{13 es primo} (Euclides):
$13\mid|A_i|$ o $13\mid|B_j|$, y como $1\le|A_i|,|B_j|\le13$, esto fuerza $|A_i|=13$ ($X$ constante) o
$|B_j|=13$ ($Y$ constante).
\begin{itemize}[leftmargin=1.4em,itemsep=1pt]
\item \textbf{$X$ no constante ($r\ge2$):} $\Rightarrow$ $Y$ debe ser constante ($q=1$).
\item \textbf{$X$ constante:} toda $Y$ es independiente de $X$.
\end{itemize}
\begin{keybox}
\textbf{Papel de la primalidad de 13:} la independencia exige $|A_i\cap B_j|=|A_i||B_j|/n$ entero; dos
variables no triviales independientes requerirían $n=a\cdot b$ con $1<a,b<n$, imposible si $n$ es primo.
Por eso sobre $\Omega_{13}$ dos variables solo pueden ser independientes si al menos una es constante.
Con $n$ compuesto (p.\,ej.\ $n=12$) sí existirían no constantes genuinamente independientes.
\end{keybox}

% ============================================================
\section*{Apéndice — Conceptos transversales}
\begin{itemize}[leftmargin=1.4em,itemsep=2pt]
\item \textbf{Etiqueta vs.\ variable:} la etiqueta identifica; la variable describe (preguntas 1, 2, 8).
\item \textbf{Partición / $\sg(X)$:} fibras de la variable; $\sg(X)$ = funciones constantes por bloque.
\item \textbf{Esperanza condicional = proyección ortogonal} = promedio por bloque = mejor predictor.
\item \textbf{Teorema de la Proyección} $\Rightarrow$ residuo ortogonal $\Rightarrow$ esperanza total,
Pitágoras, ANOVA.
\item \textbf{Ortogonalidad lo es todo:} proyecciones se suman solo si los subespacios son ortogonales
(p.\,3); no correlación $\equiv$ ortogonalidad en $L^2$ centrado (p.\,7).
\item \textbf{$R^2$ (centrado) $\le r_n^2$ (crudo)}; iguales si $\mu=0$ o ajuste perfecto (p.\,4).
\item \textbf{Norma del supremo} no viene de producto interno (falla paralelogramo) $\Rightarrow$ no sirve
para definir esperanza condicional (p.\,6).
\end{itemize}

% ============================================================
\section*{Tabla 1: Fórmulas Clave}
\vspace{0.3em}
\begin{center}
\small
\begin{tabular}{@{}p{5.5cm}p{7.5cm}@{}}
\toprule
\textbf{Concepto} & \textbf{Fórmula}\\
\midrule
Producto interno & $\langle U,V\rangle=\E[UV]=\frac1n\sum_{i=1}^n U(\omega_i)V(\omega_i)$\\
Norma en $L^2$ & $\|U\|^2=\E[U^2]=\frac1n\sum_{i=1}^n U^2(\omega_i)$\\
Esperanza condicional & $E(Y\mid\sg(X))(\omega)=\frac{1}{|B|}\sum_{\omega'\in B}Y(\omega')$, $B=$bloque de $\omega$\\
Proyección sobre recta & $P_{\sg(v)}Y=\dfrac{\langle Y,v\rangle}{\langle v,v\rangle}v$\\
Teorema de Pitágoras & $\|Y\|^2=\|\hat Y\|^2+\|Z\|^2$\\
Descomposición ANOVA & $\Var(Y)=\Var(E(Y\mid\sg(X)))+\E(\Var(Y\mid\sg(X)))$\\
Coef. determinación (crudo) & $r_n^2=\dfrac{\|E(Y\mid\sg(X))\|^2}{\|Y\|^2}=\dfrac{a}{b}$\\
Coef. determinación (centrado) & $R^2=\dfrac{\Var(\hat Y)}{\Var(Y)}=\dfrac{a-\mu^2}{b-\mu^2}$\\
$r^2$ para años & $r^2=\dfrac{SC_\text{inter}}{SC_\text{total}}=\dfrac{\sum_a n_a(\bar x_a-\bar x)^2}{\sum_i x_i^2-n\bar x^2}$\\
Proyector ortogonal & $P_M x:=\arg\min_{y\in M}\|x-y\|$, única solución con $x-Px\perp M$\\
Esperanza total (ley) & $\E[\hat Y]=\E[E(Y\mid\sg(X))]=\E[Y]$\\
Residuo ortogonal & $Z=Y-E(Y\mid\sg(X))$, $\langle Z,W\rangle=0$ para $W\in\sg(X)$\\
Igualdad del paralelogramo & $\|x+y\|^2+\|x-y\|^2=2\|x\|^2+2\|y\|^2$ (válida en Hilbert)\\
Ecuaciones normales & $X^\top(y-X\beta)=0 \Rightarrow \hat\beta=(X^\top X)^{-1}X^\top y$\\
Matriz proyectora & $P=X(X^\top X)^{-1}X^\top$, $\hat y=Py$, $e=(I-P)y$\\
\bottomrule
\end{tabular}
\end{center}

% ============================================================
\section*{Tabla 2: Símbolos y Notación}
\vspace{0.3em}
\begin{center}
\small
\begin{tabular}{@{}p{2.8cm}p{9.2cm}@{}}
\toprule
\textbf{Símbolo} & \textbf{Significado}\\
\midrule
$\Om=\Omega_n$ & Conjunto finito de etiquetas (individuos), $|\Om|=n$\\
$\mathcal P(\Om)$ & Álgebra $\sigma$ trivial (todas las partes)\\
$P_n$ & Probabilidad uniforme empírica, $P_n(\{\omega\})=1/n$\\
$\sg(X)$ & Subespacio de $L^2$ de funciones constantes por bloque de $X$\\
$L^2(\Om,\mathcal P(\Om),P_n)$ & Espacio de Hilbert de funciones en $\Om$ de cuadrado integrable\\
$\langle U,V\rangle$ & Producto interno en $L^2$: $\E[UV]$\\
$\|U\|$ & Norma de $U$ en $L^2$: $\sqrt{\E[U^2]}$\\
$E(Y\mid\sg(X))$ & Esperanza condicional de $Y$ dada la partición $\sg(X)$\\
$\hat Y$ & Predictor óptimo: $E(Y\mid\sg(X))=P_{\sg(X)}Y$\\
$Z$ & Residuo: $Z=Y-\hat Y=Y-E(Y\mid\sg(X))$\\
$P_M$ & Proyector ortogonal sobre subespacio cerrado $M$\\
$A_i=X^{-1}(\{x_i\})$ & Fibra $i$-ésima de $X$ (bloque de individuos con $X=x_i$)\\
$r^2$ & Razón de correlación: $\eta^2=\Var(E(Y\mid\sg(X)))/\Var(Y)$\\
$R^2$ & Coef. de determinación (invariante a traslaciones)\\
$r_n^2$ & Coef. basado en momentos crudos (sobrestima si $\E[Y]\neq0$)\\
$\mu=\E[Y]$ & Esperanza (valor medio) de $Y$\\
$\Var(Y)$ & Varianza de $Y$: $\E[(Y-\mu)^2]$\\
$\mathbf 1$ & Función constante $\mathbf 1(\omega)=1$ para todo $\omega$\\
$M\perp\mathcal N$ & Subespacios ortogonales: $\langle m,n\rangle=0$ para $m\in M, n\in\mathcal N$\\
$SC_\text{inter}$ & Suma de cuadrados entre bloques (variabilidad explicada)\\
$SC_\text{intra}$ & Suma de cuadrados dentro de bloques (variabilidad residual)\\
$SC_\text{total}$ & Suma total de cuadrados\\
$X\perp Y$ & Variables independientes: $P(X=x_i,Y=y_j)=P(X=x_i)P(Y=y_j)$\\
\bottomrule
\end{tabular}
\end{center}

% ============================================================
\section*{Tabla 3: Resultados Clave y Teoremas}
\vspace{0.3em}
\begin{center}
\small
\begin{tabular}{@{}p{5.0cm}p{8.0cm}@{}}
\toprule
\textbf{Teorema/Resultado} & \textbf{Enunciado}\\
\midrule
\textbf{Definición de Proyección} & $Px$ minimiza $\|x-y\|$ sobre $M\iff(x-Px)\perp M$\\
\textbf{Propiedades de $P$} & Lineal, idempotente ($P^2=P$), autoadjunto, contracción\\
\textbf{Esperanza = Proyección} & $E[X\mid\mathcal G]=P_M X$, único $\sg(\mathcal G)$-medible minimizador en $L^2$\\
\textbf{Insesgamiento} & $\E[Y-E(Y\mid\sg(X))]=0$ (como $\mathbf 1\in\sg(X)$)\\
\textbf{Ortogonalidad del residuo} & $Y-E(Y\mid\sg(X))\perp\sg(X)$ (Teorema de la Proyección)\\
\textbf{Esperanza total} & $\E[E(Y\mid\sg(X))]=\E[Y]$ (ley de la esperanza total)\\
\textbf{Teorema de Pitágoras} & $\|Y\|^2=\|\hat Y\|^2+\|Y-\hat Y\|^2$\\
\textbf{ANOVA / Descomposición} & $\Var(Y)=\underbrace{\Var(\hat Y)}_{\text{inter}}+\underbrace{\E[Var(Y\mid\sg(X))]}_{\text{intra}}$\\
\textbf{Relación $R^2$ vs $r_n^2$} & $R^2\le r_n^2$, iguales sii $\E[Y]=0$ o $Y=\hat Y$ c.s.\\
\textbf{Suma de proyecciones} & $P_{M\oplus N}=P_M+P_N$ solo si $M\perp N$ (falla sin ortogonalidad)\\
\textbf{Frisch--Waugh--Lovell} & Regresión parcial: $P_{\sg(X_1,X_2)}Y$ requiere ortogonalizar si correlacionados\\
\textbf{Coherencia jerárquica} & Particiones anidadas: $\sg(D)\subseteq\sg(L)\subseteq\sg(M)$ (Distrito $\subseteq$ Local $\subseteq$ Mesa)\\
\textbf{Falacia ecológica} & Atribuir efecto agregado al nivel individual falla (Robinson, 1950)\\
\textbf{Paradoja de Simpson} & Dirección de correlación invierte entre nivel agregado e individual\\
\textbf{Primalidad y independencia} & Con $n=13$ primo, dos variables no triviales no pueden ser independientes\\
\textbf{Igualdad del paralelogramo} & $\|x+y\|^2+\|x-y\|^2=2(\|x\|^2+\|y\|^2)$ caracteriza espacios de Hilbert\\
\textbf{Estimador MCO} & $\hat\beta=(X^\top X)^{-1}X^\top y$ es BLUE bajo homocedasticidad\\
\textbf{Ecuaciones normales} & $X^\top e=0 \Rightarrow$ residuos ortogonales a regresores\\
\bottomrule
\end{tabular}
\end{center}

\end{document}
